% Modernised reading — fonds Alexandre Grothendieck, shelfmark 40, pages 1-27.
%
% This is an INTERPRETATION, not a transcription. It re-reads the pages in
% today's notation and vocabulary, states the hypotheses the manuscript leaves
% implicit, and drops the critical apparatus so that the mathematics reads
% straight through. Everything the brackets and underlines carried moves into
% a footnote. Where the manuscript is loose or mistaken, the true statement is
% what appears here and the footnote says what the page has. In any dispute of
% substance the transcription governs: whoever cites Grothendieck must cite the
% batch-NN.fr.tex files.
%
% Pass: Opus 5.5 (claude-opus-5-5), 2026-10-09 — first pass, unchecked against
% the pages by a human. The folder was read whole: two batches, pages 1-27,
% all transcribed. This reading was made on Opus 5.5 and is not measured
% against the Opus 5 reference reading of folder 115.
%
% WHAT THE FOLDER IS. Twenty-seven pages of fast working notes, all in his
% hand, undated, no third-party material. The spine: affine group schemes as
% commutative Hopf algebras (pp. 1-2), their actions (pp. 2-4), principal
% homogeneous spaces in the affine category and two criteria for them (pp.
% 4-15), the example of the Frobenius kernel alpha_p in characteristic p
% (pp. 16-18), the classification H^1(A, alpha_p) = A/A^p (pp. 18-23), and a
% global version by the Frobenius exact sequence (pp. 24-27).
%
% NOTATION fixed once for the whole document (departures from the pages):
%   G = Spec B the group, E = Spec C the space; Delta_B for his p_B (the
%   diagonal); S for his antipode x |-> x-check; m_C for the multiplication
%   C (x) C -> C, which the pages call Delta_C (pp. 5, 10-13) — renamed, since
%   Delta also names a comultiplication; U = Hom_A(B, A) (his curled U) with
%   Delta_U and eta; C^G for the invariants, which the pages call C' (p. 6) —
%   C' is also used for a second module (p. 3) and for the dual Hom_A(C, A)
%   (p. 14), here M^vee. p_1, p_2 are kept as the pages use them: the algebra
%   maps C -> C (x) C dual to the projections. The group of pp. 16-27 is
%   named alpha_p; the pages call it B (and its sheaf the ornate B_X).
%
% Substantive departures from the pages, each footnoted where it occurs:
%   - p. 1: the counit condition « p_B x - x(x)1 - 1(x)x in J(x)J » holds for
%     x in J, not for all x (it fails for x = 1);
%   - pp. 3-4: with the convolution (uv)(b) = (u (x) v) Delta_B(b), a left
%     coaction gives a RIGHT U-module; the page fixes no order, the opposite
%     convolution is chosen so that C is a left module;
%   - pp. 9, 14-15: the converse « psi iso => phi iso » needs C finite
%     projective; the page's construction stops at « Pour ceci, ---- », the
%     proof given (psi is the C-transpose of phi) is ours;
%   - pp. 16-17: freeness of C is not a consequence of the pseudo-torsor
%     condition (p. 8's own example C = B/JB shows it);
%   - p. 20: the framed local argument needs D(C) to be a direct summand,
%     not shown; a reduction to the residue field is supplied;
%   - p. 21: the struck proof takes the FIRST nonzero coefficient; it must be
%     the last; the noetherian hypothesis queried there is not needed;
%   - p. 22: the additive law of H^1 does correspond to that of A/A^p (the
%     page says it is not proved); the corollary must keep u^p in A
%     (otherwise A[u]/(u^p - u) would qualify);
%   - p. 23: the classes au + b fill only part of P(C/A), all of it (minus 0)
%     over a field;
%   - pp. 24-25: the exact sequence is one of fppf sheaves, and Frobenius is
%     not surjective in the Zariski or étale sense;
%   - p. 26: n ranges over all positive integers prime to p, not 1..p-1.
%
% Announced and never established in the folder: the converse of pp. 14-15
% (left at « Pour ceci, ---- » with « ??? »); the two « Questions » of p. 18
% (only the second is pursued); the additive law of p. 22; the surjectivity
% of Frobenius « en un sens technique » (p. 25, margin « ? »).

\documentclass[11pt,a4paper]{article}
\input{../preamble/grothendieck}

\folder{40}
\pages{1}{27}
\dating{[années 1960-1970]}
\watermark{Édition de démonstration\\interprétation personnelle de l'œuvre}
\foldertitle{Calculs [affines sur des] fibrés principaux affines : notes manuscrites (s.d.) — lecture modernisée du dossier entier}

\begin{document}

\section*{Résumé}

\begin{resume}
Ces vingt-sept pages sont des calculs de travail, écrits vite, sur une
question simple à dire : quand un groupe agit sur un objet de façon
« parfaitement régulière », que peut-on dire de cet objet, et combien y en
a-t-il ?

Le modèle à garder en tête est celui d'un groupe $G$ qui agit sur un ensemble
$E$ \emph{librement et transitivement} : pour deux points $x$, $y$ de $E$, il
existe un et un seul $g$ tel que $x = g\,y$. Un tel $E$ ressemble trait pour
trait à $G$ lui-même, à ceci près qu'on n'y a pas choisi d'origine --- comme
un espace affine ressemble à un espace vectoriel dont on aurait oublié où est
le zéro. On l'appelle un \emph{espace principal homogène}, ou aujourd'hui un
\emph{torseur}. Sur un ensemble il n'y a rien à classer : dès qu'on choisit un
point, $E$ s'identifie à $G$. Mais si $E$ dépend d'un paramètre, on ne peut
pas toujours choisir ce point de façon cohérente, et les torseurs non
triviaux mesurent exactement cette impossibilité.

Ici, groupes et espaces ne sont pas des ensembles mais des objets de la
géométrie algébrique, connus par leurs algèbres de fonctions. Le dossier
commence donc par traduire tout le langage des groupes dans celui des
algèbres : la loi de groupe devient une « diagonale » $B \to B \otimes B$,
l'action devient une application $C \to B \otimes C$, et la condition « libre
et transitive » devient l'affirmation qu'une certaine application
$C \otimes C \to B \otimes C$ est un isomorphisme. Quand l'algèbre $B$ est de
dimension finie, on peut aussi passer à son dual, qui est une algèbre
d'\emph{opérateurs} agissant sur $C$, et la condition se lit sur ces
opérateurs : ils doivent fournir une base de tous les endomorphismes de $C$.

Puis vient un exemple qui n'a pas d'équivalent pour les ensembles. En
caractéristique $p$, le groupe des $x$ tels que $x^p = 0$ n'a qu'un seul
point, et pourtant n'est pas trivial : il est « infinitésimal ». Agir par ce
groupe, c'est la même chose que se donner une dérivation $D$ de puissance
$p$-ième nulle. Le dossier montre que ses torseurs sont exactement les
algèbres obtenues en adjoignant une racine $p$-ième, $A[U]/(U^p - a)$, et que
deux d'entre elles sont isomorphes si et seulement si les nombres $a$
diffèrent d'une puissance $p$-ième. La classification tient en une formule :
le groupe de ces torseurs est $A/A^p$.

Les dernières pages refont ce résultat d'un seul coup d'œil, à l'aide d'une
suite exacte --- le groupe infinitésimal est le noyau de l'élévation à la
puissance $p$ --- et en tirent une version globale, valable sur une variété
quelconque. Une remarque finale montre un phénomène inattendu : deux torseurs
« lisses » peuvent avoir une somme qui ne l'est pas.

Ce qu'on cherchera ensuite : schémas en groupes affines et algèbres de Hopf,
dualité de Cartier, torseurs et cohomologie plate, le schéma en groupes
$\alpha_p$, et les extensions radicielles de hauteur un.

\keywords{affine group scheme, Hopf algebra, comodule algebra, module algebra, pseudo-torsor, torsor, Cartier duality, Frobenius kernel, restricted Lie algebra, purely inseparable extension, fppf cohomology}
\end{resume}

\section*{Le fil du dossier, et les conventions}

Le dossier a une seule ligne, sans reprise ni mouture.

\begin{itemize}
\item \textbf{Pages 1 et 2} --- les schémas en groupes affines comme
algèbres de Hopf commutatives ; le dual $\mathcal{U}$ quand l'algèbre est
libre de rang fini.
\item \textbf{Pages 2 à 4} --- les actions, sous forme de coaction, puis sous
forme de module sur $\mathcal{U}$.
\item \textbf{Pages 4 à 10} --- les espaces principaux homogènes : la
condition d'isomorphisme, le quotient par le groupe, et deux critères
équivalents pour $B$ et $C$ libres de type fini.
\item \textbf{Pages 11 à 15} --- reprise catégorique et démonstration de
l'équivalence des deux critères, dont une direction reste inachevée sur la
page.
\item \textbf{Pages 16 à 18} --- l'exemple de $\alpha_p$ : actions par
dérivations, et la « Proposition ».
\item \textbf{Pages 18 à 23} --- la classification : $H^1(A,\alpha_p) = A/A^p$,
le « Théorème » et son « Corollaire ».
\item \textbf{Pages 24 à 27} --- la suite exacte de Frobenius, la version
globale, et la « Remarque » sur $k[[X]]$.
\end{itemize}

Conventions valables pour tout le dossier. $A$ est un anneau commutatif, les
algèbres sont commutatives et unitaires sauf mention contraire, et les produits
tensoriels sont pris sur $A$. Le groupe est $G = \operatorname{Spec} B$,
l'espace sur lequel il agit est $E = \operatorname{Spec} C$. On note
$\Delta_B$ la comultiplication de $B$ (la page écrit $p_B$), $\varepsilon$ sa
coünité, $J = \operatorname{Ker}\varepsilon$, $S$ son antipode (la page écrit
$x \mapsto \check{x}$), et $m_C \colon C \otimes C \to C$ la multiplication
de $C$.\footnote{La page note $\Delta_C$ la multiplication de $C$ (pages 5,
10 à 13). On la renomme $m_C$ parce que $\Delta$ désigne ailleurs une
comultiplication, et que les deux se rencontrent dans les mêmes formules.}
On note $p_1, p_2 \colon C \to C \otimes C$ les homomorphismes
$x \mapsto x \otimes 1$ et $x \mapsto 1 \otimes x$, duaux des deux projections
$E \times E \to E$, comme le fait la page. Les invariants de $C$ sont notés
$C^G$.\footnote{La page les note $C'$ (page 6). Elle emploie aussi $C'$ pour
un second module quelconque (page 3) et pour le dual $\operatorname{Hom}_A(C,
A)$ (page 14) ; ce dernier est noté ici $C^{\vee}$.} À partir de la page 16
le groupe est $\alpha_p$, que la page appelle $B$, et son faisceau sur un
préschéma $X$ est $\alpha_{p,X}$.

\pagerange{1}{2}
\section*{Groupes affines et algèbres de Hopf (pages 1 et 2)}

Un groupe dans la catégorie opposée à celle des $A$-algèbres --- c'est-à-dire
un schéma en groupes affine $G = \operatorname{Spec} B$ sur $A$ --- est une
$A$-algèbre $B$ munie d'un homomorphisme d'algèbres coassociatif
$\Delta_B \colon B \to B \otimes B$, d'une coünité
$\varepsilon \colon B \to A$ et d'une antipode $S \colon B \to B$ : une
\emph{algèbre de Hopf commutative}. La coünité est caractérisée par
\[
\Delta_B x - (x \otimes 1 + 1 \otimes x) \in J \otimes J
\qquad (x \in J),
\]
et elle est unique si elle existe.\footnote{La page écrit cette condition pour
tout $x \in B$. Elle est fausse pour $x = 1$, puisque
$\Delta_B(1) - 2\,(1 \otimes 1) = -\,1 \otimes 1 \notin J \otimes J$ ; elle
est juste pour $x \in J$, ou, pour tout $x$, sous la forme
$\Delta_B x - x \otimes 1 - 1 \otimes x + \varepsilon(x)\, 1 \otimes 1 \in J
\otimes J$. Comme $B = A \cdot 1 \oplus J$, la condition sur $J$ équivaut aux
deux axiomes de coünité $(\varepsilon \otimes \mathrm{id})\Delta_B =
(\mathrm{id} \otimes \varepsilon)\Delta_B = \mathrm{id}$.} L'antipode est
l'homomorphisme d'algèbres tel que, si $\Delta_B x = \sum y_i \otimes z_i$,
\[
\sum S(y_i)\, z_i = \sum y_i\, S(z_i) = \varepsilon(x)\cdot 1 .
\]
Elle est unique, c'est une involution (parce que $B$ est commutative), et
l'on a $\varepsilon \circ S = \varepsilon$ et
$\Delta_B (S x) = \sum S(z_i) \otimes S(y_i)$.

Supposons $B$ libre de rang fini sur $A$\footnote{La page écrit « libre de
rang fini sur $B$ » ; il faut lire $A$.} et soit
$\mathcal{U} = \operatorname{Hom}_A(B, A)$ son dual. Toute la structure passe
au dual en se retournant : la multiplication de $B$ devient une
comultiplication $\Delta_{\mathcal{U}}$, \emph{cocommutative} et
coassociative, avec la coünité $\eta(u) = u(1_B)$ ; la comultiplication de
$B$ devient une multiplication associative et unitaire, en général non
commutative ; les deux sont compatibles ($\Delta_{\mathcal{U}}$ est un
homomorphisme d'algèbres) ; l'antipode devient un antiautomorphisme
involutif $S_{\mathcal{U}}$, déterminé de façon unique par l'identité
$\sum y_i\, S_{\mathcal{U}}(z_i) = \sum S_{\mathcal{U}}(y_i)\, z_i =
\eta(x)\cdot 1$ pour $\Delta_{\mathcal{U}} x = \sum y_i \otimes z_i$. Bref,
$\mathcal{U}$ est une algèbre de Hopf cocommutative, l'« hyperalgèbre » du
groupe.\footnote{Le mot est celui de Dieudonné, qui l'introduit au cours des
années cinquante pour les groupes de Lie formels en caractéristique $p$. On
dit aujourd'hui algèbre des distributions (au sens de Demazure--Gabriel) ou,
pour un groupe fini et plat comme ici, algèbre du groupe ; quand $G$ est
commutatif, $\mathcal{U}$ est commutative et son spectre est le \emph{dual de
Cartier} de $G$. Le dossier emploie le mot plus librement, pour toute algèbre
de Hopf (page 24) ; la page ne dit pas où elle l'a pris.}

\pagerange{2}{4}
\section*{Actions : coactions et modules-algèbres (pages 2 à 4)}

Faire agir $G$ à gauche sur $E = \operatorname{Spec} C$, c'est se donner un
homomorphisme d'algèbres unitaires
\[
q_C \colon C \longrightarrow B \otimes C
\]
qui satisfait à deux conditions : l'\emph{associativité},
$(\mathrm{id}_B \otimes q_C)\, q_C = (\Delta_B \otimes \mathrm{id}_C)\,
q_C$, et l'\emph{unité}, $(\varepsilon \otimes \mathrm{id}_C)\, q_C =
\mathrm{id}_C$, qui s'écrit aussi $q_C(y) \in 1 \otimes y + J \otimes C$.
C'est ce qu'on appelle aujourd'hui une \emph{algèbre-comodule} sur $B$.

\begin{tikzcd}[column sep=huge, row sep=large]
B \otimes B \otimes C & B \otimes C \arrow[l, "\mathrm{id}_B \otimes q_C"'] \\
B \otimes C \arrow[u, "\Delta_B \otimes \mathrm{id}_C"] & C \arrow[l, "q_C"] \arrow[u, "q_C"']
\end{tikzcd}

Si $B$ est libre de rang fini, $B \otimes C' \simeq
\operatorname{Hom}_A(\mathcal{U}, C')$ pour tout $A$-module $C'$, de sorte
qu'une application $A$-linéaire $C \to B \otimes C'$ est la même chose qu'une
application $A$-linéaire $\mathcal{U} \to \operatorname{Hom}_A(C, C')$. Pour
$C' = C$, la coaction devient ainsi une action de $\mathcal{U}$ sur $C$, que
l'on note $u \cdot x = (u \otimes \mathrm{id}_C)\, q_C(x)$. Les axiomes
d'associativité et d'unité disent que $C$ est un $\mathcal{U}$-module
unitaire\footnote{Un point d'ordre que la page ne fixe pas. Si l'on munit
$\mathcal{U}$ de la convolution $(uv)(b) = (u \otimes v)\,\Delta_B(b)$, une
coaction \emph{à gauche} fait de $C$ un module \emph{à droite} : on trouve
$(uv)\cdot x = v \cdot (u \cdot x)$. Pour que $C$ soit un module à gauche,
comme la page l'écrit, on prend la convolution opposée,
$(uv)(b) = (v \otimes u)\,\Delta_B(b)$, et c'est la convention adoptée ici.
Quand $G$ est commutatif --- c'est le cas de tous les exemples du dossier ---
les deux coïncident.} ; et que $q_C$ soit un homomorphisme d'algèbres
unitaires se traduit par
\[
u \cdot 1_C = \eta(u)\, 1_C, \qquad
u \cdot (xy) = \sum (v_i \cdot x)(w_i \cdot y)
\quad\text{si } \Delta_{\mathcal{U}} u = \sum v_i \otimes w_i .
\]
C'est la notion de \emph{$\mathcal{U}$-module-algèbre}.\footnote{Dans la
première de ces conditions la page écrit $q_{\mathcal{U}}$, dans la seconde
$p_{\mathcal{U}}$, pour le même homomorphisme $\mathcal{U} \to
\operatorname{Hom}_A(C, C')$ ; on n'en garde qu'un. La page écrit d'abord ces
conditions pour deux algèbres $C$ et $C'$ distinctes, puis pour $C' = C$.}
Si $u$ est un élément « primitif » ($\Delta_{\mathcal{U}} u = u \otimes 1 +
1 \otimes u$), la seconde condition dit simplement que $u$ agit par une
dérivation : c'est ce qui servira à partir de la page 16.

\pagerange{4}{8}
\section*{Espaces principaux homogènes et quotient (pages 4 à 8)}

\subsection*{La condition}

Que $E$ soit un espace principal homogène sous $G$ --- que l'action soit
« libre et transitive » --- s'exprime en demandant que le morphisme
$G \times E \to E \times E$, $(g, y) \mapsto (g y, y)$, soit un isomorphisme.
Sur les algèbres, c'est l'homomorphisme
\[
\varphi = (q_C, p_2) \colon C \otimes C \longrightarrow B \otimes C,
\qquad
\varphi(x \otimes y) = q_C(x)\,(1 \otimes y),
\]
c'est-à-dire le composé $(\mathrm{id}_B \otimes m_C) \circ (q_C \otimes
\mathrm{id}_C)$. Il revient au même de demander qu'il existe un
homomorphisme d'algèbres unitaires $g_C \colon B \to C \otimes C$ --- le
morphisme $E \times E \to G$ qui envoie $(x, y)$ sur l'unique $g$ tel que
$x = g y$ --- rendant commutatifs les deux triangles
\[
(g_C, p_2) \circ q_C = p_1, \qquad \varphi \circ g_C = p_1,
\]
où $(g_C, p_2) \colon B \otimes C \to C \otimes C$ est
$b \otimes y \mapsto g_C(b)\,(1 \otimes y)$ et où $p_1$ désigne, dans le
second, l'homomorphisme $b \mapsto b \otimes 1$ de $B$ dans
$B \otimes C$. En effet $(g_C, p_2)$ est alors l'inverse de
$\varphi$.\footnote{Le haut de la page 5, où ces conditions sont d'abord
cherchées, est barré et à peu près illisible ; la transcription en donne une
lecture très incomplète, qu'on ne reconstitue pas. La formulation retenue est
celle que la page 11 reprend en termes catégoriques.}

Cette condition seule ne fait pas de $E$ un torseur au sens d'aujourd'hui :
elle est satisfaite, par exemple, par l'objet vide ($C = 0$). On dit qu'un
tel $E$ est \emph{formellement principal homogène}, ou un
\emph{pseudo-torseur}.\footnote{Le terme « formellement principal homogène »
est celui de la page 11 ; c'est aussi celui de SGA 3. Un \emph{torseur}
demande en outre que $E$ soit, localement pour une topologie convenable (la
topologie plate), isomorphe à $G$ --- pour $E$ affine, que $C$ soit
fidèlement plate sur $A$. La page 10 y ajoute une condition sur les
invariants ; quand $C$ est libre et non nul, comme la page le suppose à
partir de la page 9, $C$ est fidèlement plate et les deux notions
coïncident.}

\subsection*{Le quotient par le groupe}

Soit $h \colon \mathcal{D} \to C$ un homomorphisme d'algèbres tel que
$q_C \circ h = (x \mapsto 1 \otimes x) \circ h$, autrement dit un morphisme
$E \to \operatorname{Spec}\mathcal{D}$ constant sur les orbites. Cela
signifie exactement que $h$ prend ses valeurs dans la sous-algèbre des
\emph{invariants}
\[
C^G = \{\, x \in C \;:\; q_C(x) = 1 \otimes x \,\}.
\]
Donc \emph{le quotient de $E$ par $G$ dans la catégorie des schémas affines
existe, et c'est $\operatorname{Spec} C^G$} : c'est l'égalisateur des deux
homomorphismes $C \rightrightarrows B \otimes C$.\footnote{La page appelle
ces éléments « primitifs », puis ajoute au-dessus « invariants », qui est le
mot juste ; à la page 7 elle dit « semi-primitifs ». La lecture d'une ligne de
la page 6 (« on voit qu'on a tout ce qu'il faut ») est très incertaine.} La
question que pose la page 6 --- un tel $h$ se factorise-t-il par
$A$ ? --- revient donc à savoir si $C^G = A$.

Mais ce quotient « dans la catégorie » peut être très différent du quotient
géométrique. Soit $k$ algébriquement clos, $E$ le groupe linéaire, $G$ un
sous-groupe de Borel agissant par translations. Le quotient $E/G$ existe dans
la catégorie des schémas : c'est la variété des drapeaux, qui est complète.
Ses seules fonctions globales sont les constantes ; donc
$C^G = k$, et le quotient dans la catégorie des variétés affines est réduit à
un point.\footnote{Paragraphe embrassé d'un grand crochet en marge, à la
lecture incertaine en plusieurs mots (« groupe linéaire », « réduit au
point »).}

\subsection*{Les invariants d'un pseudo-torseur}

Si $E$ est un pseudo-torseur et $x \in C^G$, le premier triangle appliqué à
$x$ donne
\[
x \otimes 1 = 1 \otimes x \quad\text{dans } C \otimes C .
\]
Si de plus $A \cdot 1_C$ est un facteur direct du $A$-module $C$, avec
$\lambda \mapsto \lambda 1_C$ injectif, cette égalité force
$x \in A \cdot 1_C$ : écrivant $C = A \cdot 1 \oplus N$ et $x = \lambda + n$,
la composante de $x \otimes 1 - 1 \otimes x$ dans $N \otimes A$ est
$n \otimes 1$, d'où $n = 0$. Donc $C^G = A$. L'hypothèse est satisfaite dès
que $C$ est un $A$-module libre ayant une base contenant $1_C$, et en
particulier dès que $C$ est libre de type fini et non nul.\footnote{La page
dit « comme on voit facilement ». Pour $C$ projectif de type fini, un
élément est unimodulaire dès que son image dans chaque fibre
$C \otimes k(\mathfrak{p})$ est non nulle ; c'est le cas de $1_C$ dès que
toutes ces fibres sont non nulles. L'hypothèse $C \neq 0$, que la page ne
dit pas, est nécessaire.}

Sans hypothèse de cette sorte, c'est faux, et la page 8 le montre : si $J$
est un idéal de $A$ et $C = B / JB$, avec la coaction obtenue en réduisant
$\Delta_B$ modulo $J$, les axiomes restent vérifiés, $E$ est même un
pseudo-torseur, et pourtant $C^G = A/J$.\footnote{Que cet $E$ soit un
pseudo-torseur n'est pas dit par la page, qui écrit seulement que « les
axiomes restent vérifiés » ; c'est immédiat, puisque $\varphi$ est la
réduction modulo $J$ de l'isomorphisme $(\Delta_B, p_2)$ de $G \times G$.
Le même exemple montre, contre ce que suggère la page 16, que la liberté de
$C$ n'est pas conséquence de la seule condition de pseudo-torseur.}

\pagerange{9}{10}
\section*{Deux critères quand $B$ et $C$ sont libres de type fini (pages 9 et 10)}

Supposons $B$ et $C$ libres de type fini, $C \neq 0$. Le paragraphe
précédent montre alors que la condition sur les invariants est conséquence de
celle de pseudo-torseur. Le critère s'écrit de deux façons.

\emph{Par $\varphi$.} Via $B \otimes C \simeq \operatorname{Hom}_A(\mathcal{U},
C)$, l'homomorphisme $\varphi$ devient
\[
\varphi \colon C \otimes C \longrightarrow \operatorname{Hom}_A(\mathcal{U}, C),
\qquad \varphi(x \otimes y)(u) = (u \cdot x)\, y ,
\]
le second membre étant muni de la convolution $v * w = m_C \circ (v \otimes
w) \circ \Delta_{\mathcal{U}}$, pour laquelle $\varphi$ est un homomorphisme
d'algèbres. $E$ est un pseudo-torseur si et seulement si $\varphi$ est
bijectif.

\emph{Par les opérateurs.} Soit
\[
\psi \colon C \otimes \mathcal{U} \longrightarrow \operatorname{End}_A(C),
\qquad \psi(x \otimes u)(y) = x\,(u \cdot y).
\]
\emph{$\varphi$ est un isomorphisme si et seulement si $\psi$ en est un} :
autrement dit, $E$ est un pseudo-torseur si et seulement si les opérateurs
de $\mathcal{U}$, pris avec des coefficients dans $C$, fournissent
exactement tous les endomorphismes $A$-linéaires de $C$.\footnote{En marge,
un « ? » : sur la page, l'équivalence est annoncée, pas démontrée. Elle l'est
aux pages 12 à 15, dans un sens seulement ; voir la section suivante. Dans la
conclusion entre crochets, la page écrit « fibré principal homogène sous
$C$ » ; il faut lire $B$. La moitié inférieure de la page 9, barrée, esquisse
un argument par changement de base $A' = C$ qu'on ne reconstitue pas.}

Enfin, la condition sur les invariants s'écrit en termes de $\mathcal{U}$ :
$x \in C^G$ si et seulement si $u \cdot x = \eta(u)\, x$ pour tout
$u \in \mathcal{U}$.

\pagerange{11}{15}
\section*{Reprise catégorique ; l'équivalence des deux critères (pages 11 à 15)}

\subsection*{Le cadre}

Soit $\mathcal{C}$ une catégorie avec objet final $e$ et produits, $G$ un
groupe de $\mathcal{C}$, $E$ un objet sur lequel $G$ opère par
$q_E \colon G \times E \to E$. On dit que $E$ est \emph{formellement
principal homogène} si $(q_E, p_2) \colon G \times E \to E \times E$ est un
isomorphisme ; il revient au même qu'il existe $g_E \colon E \times E \to G$
--- « $x = g(x, y)\, y$ », dit la marge --- tel que
\[
q_E \circ (g_E, p_2) = p_1 \ \text{sur } E \times E,
\qquad
g_E \circ (q_E, p_2) = p_1 \ \text{sur } G \times E .
\]
Pour $\mathcal{C}$ opposée à la catégorie des $A$-algèbres, on retrouve
exactement les deux triangles des pages 4 et 5.

\subsection*{Le critère $\psi$ est le transposé du critère $\varphi$}

On n'a plus besoin ici que $B$ soit libre. Posons
\[
\psi \colon \operatorname{Hom}_A(B, C) \longrightarrow \operatorname{End}_A(C),
\qquad
\psi(v) = m_C \circ (v \otimes \mathrm{id}_C) \circ q_C .
\]
Quand $B$ est libre de type fini, $\operatorname{Hom}_A(B, C) \simeq C
\otimes \mathcal{U}$ et l'on retrouve le $\psi$ de la page 9.

\emph{Si $\varphi$ est un isomorphisme, $\psi$ en est un}, d'inverse
\[
\psi'(w) = m_C \circ (w \otimes \mathrm{id}_C) \circ g_C .
\]
La page le vérifie par deux colonnes de diagrammes (page 13), qui réécrivent
$\psi\psi'(w)$ et $\psi'\psi(v)$ en échangeant l'ordre des étapes jusqu'à
reconnaître $w$ et $v$.\footnote{La page donne, pour chacune des deux
identités, six colonnes reliées par des accolades ; la transcription ne
reproduit que la première et la dernière. Dans les colonnes intermédiaires de
la seconde, plusieurs $B$ figurent là où l'on attend $C$ ; ce sont des
lapsus sans incidence sur le calcul.} Il y a une raison plus courte, qui
explique aussi la réciproque. Munissons $B \otimes C$ et $C \otimes C$ de
leur structure de $C$-module par le facteur de droite ; $\varphi$ est alors
$C$-linéaire. Une application $A$-linéaire $v \colon B \to C$ est la même
chose qu'une application $C$-linéaire $\tilde v \colon B \otimes C \to C$,
et de même $\operatorname{End}_A(C) \simeq \operatorname{Hom}_C(C \otimes C,
C)$, $w \mapsto (x \otimes y \mapsto w(x)\,y)$. Dans ces identifications,
\[
\psi(v) = \tilde v \circ \varphi \ \text{restreint à } C \otimes 1 :
\qquad \psi = \operatorname{Hom}_C(\varphi, C).
\]
Le critère $\psi$ est donc le \emph{transposé}, en tant que $C$-module, du
critère $\varphi$.\footnote{Cette lecture est la nôtre ; la page procède par
calcul direct.} Il en résulte aussitôt que $\psi$ est bijectif si $\varphi$
l'est.

\subsection*{La réciproque}

Réciproquement, \emph{si $C$ est un $A$-module projectif de type fini (par
exemple libre de type fini) et $B$ libre de type fini, et si $\psi$ est un
isomorphisme, alors $\varphi$ en est un}. En effet $B \otimes C$ et
$C \otimes C$ sont alors des $C$-modules projectifs de type fini, pour
lesquels un homomorphisme est bijectif dès que son transposé l'est.

La page suit un autre chemin, qu'elle ne mène pas à son terme. Elle part de
l'inverse $\psi'$ de $\psi$, le regarde comme une application
$\psi'_1 \colon B \to \operatorname{Hom}(\operatorname{End}_A(C), C)$, et
cherche à en tirer $g_C$ en composant avec une application naturelle
\[
\rho \colon \operatorname{Hom}(\operatorname{End}_A(C), C) \simeq C \otimes
C^{\vee} \otimes C \xrightarrow{\ \mathrm{id} \otimes \varepsilon_C \otimes
\mathrm{id}\ } C \otimes C,
\]
où $\varepsilon_C \colon C^{\vee} \to A$ est l'évaluation en $1_C$. Il
resterait à vérifier que $g_C = \rho \circ \psi'_1$ satisfait aux deux
triangles ; la page s'arrête sur « Pour ceci, ---- », avec « ??? » en marge,
et le reste de la page 15 est blanc.\footnote{La page 14 contient aussi un
crochet barré qui énonce l'application naturelle
$P \otimes \operatorname{Hom}(Q, R) \to
\operatorname{Hom}(\operatorname{Hom}(P, Q), R)$. L'argument par
transposition donné ci-dessus est le nôtre ; il établit la réciproque sans
passer par $\rho$, et ne dit pas si le $g_C$ de la page est le bon.}

\pagerange{16}{18}
\section*{L'exemple de $\alpha_p$ (pages 16 à 18)}

\subsection*{Le groupe et son hyperalgèbre}

Soit $A$ de caractéristique $p$ et
\[
\mathcal{U} = A[X]/(X^p),
\qquad \Delta X = X \otimes 1 + 1 \otimes X,
\quad \varepsilon(X) = 0,
\quad S(X) = -X .
\]
La comultiplication est bien définie parce que, en caractéristique $p$,
$\Delta(X^p) = X^p \otimes 1 + 1 \otimes X^p$ appartient à l'idéal engendré
par $X^p \otimes 1$ et $1 \otimes X^p$. C'est l'algèbre de Hopf du schéma en
groupes $\alpha_p$ des $x$ tels que $x^p = 0$ --- le noyau du Frobenius du
groupe additif, comme le diront les pages 24 et suivantes. Elle est isomorphe
à sa duale : la base duale de $(X^i)_{0 \le i < p}$ est formée des
$Y^i / i!$, où $Y$ est primitif et $Y^p = 0$.\footnote{La page affirme cet
autodualité sans la démontrer ; la vérification est la nôtre. Elle utilise
que $i!$ est inversible pour $i < p$. En termes modernes, $\alpha_p$ est son
propre dual de Cartier.} On peut donc regarder $\mathcal{U}$ indifféremment
comme l'algèbre du groupe ou comme son hyperalgèbre, et la page passe de l'un
à l'autre : elle appelle $B$ le groupe dont $\mathcal{U}$ est le dual.

\subsection*{Actions par dérivations}

Faire opérer $\mathcal{U}$ sur une $A$-algèbre commutative $C$, c'est se
donner l'image $D$ de $X$, un endomorphisme $A$-linéaire de $C$, soumis à deux
conditions : $D^p = 0$, pour que $A[X]/(X^p)$ agisse ; et $D(xy) = x\,Dy +
y\,Dx$, parce que $X$ est primitif. \emph{Une action de $\alpha_p$ sur
$\operatorname{Spec} C$ est donc la même chose qu'une $A$-dérivation de $C$
de puissance $p$-ième nulle.}\footnote{En caractéristique $p$, $D^p$ est
encore une dérivation ; la condition $D^p = 0$ n'est donc pas automatique,
et c'est elle qui distingue $\alpha_p$ du groupe additif tout entier.}

Supposons $C$ libre de type fini et non nul. Comme $C \otimes C \simeq B
\otimes C$, son rang est nécessairement $p$.\footnote{La page suppose $C$
libre et ajoute que cela « semble d'ailleurs une conséquence de notre
hypothèse ». Ce n'est pas le cas pour un simple pseudo-torseur (exemple de la
page 8) ; c'est vrai, localement, pour un torseur : si $C$ est fidèlement
plate, l'isomorphisme $C \otimes C \simeq B \otimes C$ montre par descente
que $C$ est localement libre de rang $p$.} Par le critère de la page 9,
$\operatorname{Spec} C$ est alors principal homogène si et seulement si
\[
1, D, D^2, \ldots, D^{p-1} \ \text{forment une base du $C$-module}\
\operatorname{End}_A(C),
\]
où $C$ agit sur $\operatorname{End}_A(C)$ par $(\lambda w)(y) = \lambda\,
w(y)$.

\subsection*{La « Proposition »}

Supposons que $C$ admette une $p$-base sur $A$ réduite à un élément $u$ :
$C \simeq A[u]/(u^p - a)$ avec $a \in A$. Toute valeur de $Du$ définit une
$A$-dérivation, puisque $D(u^p - a) = p\, u^{p-1} Du = 0$ ; prenons $Du = 1$,
de sorte que $D(u^n) = n\, u^{n-1}$ et $D^p = 0$. Identifions
$\operatorname{End}_A(C)$ à $C^p$ par $w \mapsto (w(u^j))_{0 \le j \le p-1}$
--- c'est un isomorphisme de $C$-modules. Les $D^i$ deviennent les lignes de
la matrice triangulaire
\[
\bigl(D^i(u^j)\bigr)_{0 \le i, j \le p-1}
=
\begin{pmatrix}
1 & u & u^2 & \cdots & u^{p-1} \\
0 & 1 & 2u & \cdots & (p-1)\, u^{p-2} \\
0 & 0 & 2! & \cdots & (p-1)(p-2)\, u^{p-3} \\
\vdots & & & \ddots & \vdots \\
0 & 0 & 0 & \cdots & (p-1)!
\end{pmatrix},
\]
de déterminant $0!\,1!\,2!\cdots(p-1)!$, image d'un élément non nul de
$\mathbb{Z}/p\mathbb{Z}$, donc inversible. D'où :

\begin{quote}
\textbf{Proposition} (page 18). \emph{Soit $C$ une $A$-algèbre admettant une
$p$-base réduite à un élément $u$, c'est-à-dire $C = A[u]/(u^p - a)$. Muni de
la dérivation $D$ telle que $Du = 1$, $\operatorname{Spec} C$ est un torseur
sous $\alpha_p$.}
\end{quote}

C'est bien un torseur et non seulement un pseudo-torseur, puisque $C$ est
libre et non nul, donc fidèlement plat.\footnote{La page dit « espace homogène
fibré principal » et n'a pas la distinction ; on la fait parce que la
condition des invariants (page 10) est ici satisfaite, $1_C$ faisant partie
d'une base.}

\pagerange{18}{18}
\subsection*{Deux « Questions »}

La page 18 pose deux questions qu'elle ne reprend pas. La première :
pour $C$ comme ci-dessus, déterminer toutes les façons d'obtenir la
$C$-algèbre $\operatorname{End}_A(C)$ comme $C \otimes \mathcal{U}$ à partir
d'une sous-$A$-hyperalgèbre $\mathcal{U}$ libre de rang $p$. La seconde :
déterminer celles de ces $\mathcal{U}$ qui sont l'algèbre enveloppante d'une
$p$-algèbre de Lie libre de rang un, de puissance $p$-ième
nulle.\footnote{Les ajouts interlinéaires de ces deux lignes sont serrés et
en partie illisibles, et l'ordre de lecture est incertain ; on ne garde que
ce que la transcription lit. Les noms modernes sont ceux d'\emph{algèbre de
Lie restreinte} (Jacobson, 1941) et d'\emph{algèbre enveloppante restreinte}.
L'équivalence entre schémas en groupes finis de hauteur un et $p$-algèbres de
Lie est exposée dans SGA 3 (exposé VII$_{\mathrm{A}}$) ; la page ne cite
rien.} Les pages 22 et 23 répondent en partie à la première, pour
$\mathcal{U} \simeq \alpha_p$.

\pagerange{19}{23}
\section*{La classification : $H^1(A, \alpha_p) = A/A^p$ (pages 19 à 23)}

\subsection*{Les invariants d'un torseur sous $\alpha_p$}

Soit $C$ une $A$-algèbre libre de rang $p$, munie d'une $A$-dérivation $D$
telle que $D^p = 0$ et que $1, D, \ldots, D^{p-1}$ forment une base de
$\operatorname{End}_A(C)$ sur $C$. Alors \emph{$\operatorname{Ker} D = A$}.

Cela repose sur un lemme d'orthogonalité. Pour un sous-$A$-module $M$ de $C$,
notons $M^{\circ} \subset \operatorname{End}_A(C)$ l'ensemble des
endomorphismes nuls sur $M$, et pour $L \subset \operatorname{End}_A(C)$
notons $L^{\circ} \subset C$ l'ensemble des éléments annulés par $L$.

\begin{quote}
\textbf{Lemme} (page 19). \emph{Si $M$ est libre de rang $m$ et facteur
direct dans $C$, alors $M^{\circ}$ est un sous-$C$-module libre de rang $p - m$
et facteur direct de $\operatorname{End}_A(C)$, et $M = (M^{\circ})^{\circ}$.
Si $L \subset M^{\circ}$ est un sous-$C$-module libre de rang $p - m$ et
facteur direct, alors $L = M^{\circ}$, et donc $M = L^{\circ}$.}
\end{quote}

Si $C = M \oplus N$, on a $M^{\circ} \simeq \operatorname{Hom}_A(N, C)$, libre
de rang $p - m$ sur $C$ ; un $x \notin M$ a une composante non nulle dans
$N$, que la projection sur $N$ ne tue pas ; et $L$, facteur direct de
$M^{\circ}$ de même rang, lui est égal. On applique le lemme à $M = A \cdot
1_C$ et à $L$ engendré par $D, D^2, \ldots, D^{p-1}$, qui sont nuls sur $A$ et
font partie d'une base : il vient
$A = L^{\circ} = \{x : D^i x = 0,\ 1 \le i \le p-1\} = \operatorname{Ker}
D$.\footnote{L'énoncé du lemme est en partie illisible sur la page (une
dizaine de mots), et deux notes tournées dans la marge (dont « Lemme ») ne
se lisent qu'à peine. La formulation ci-dessus est reconstituée à partir de
ce qui est lisible et de l'usage qui en est fait à la page 20 ; la
démonstration esquissée est la nôtre.}

Une conséquence au passage : pour tout $x \in C$, $x^p \in A$, puisque
$D(x^p) = p\,x^{p-1}\,Dx = 0$.

\subsection*{L'élément $u$}

\begin{quote}
\textbf{Corollaire} (page 20). \emph{Il existe $u \in C$ tel que $Du = 1$,
et $u$ est unique à l'addition près d'un élément de $A$.}
\end{quote}

L'unicité vient de $\operatorname{Ker} D = A$. Pour l'existence, la question
est locale sur $\operatorname{Spec} A$ : l'ensemble des $u$ tels que
$Du = 1$, s'il n'est pas vide, est un espace principal homogène sous le
groupe additif $A$, et un torseur sous le faisceau structural d'un schéma
affine est trivial ($H^1(\operatorname{Spec} A, \mathcal{O}) = 0$) ; on peut
donc supposer $A$ local, d'idéal maximal $\mathfrak{m}$ et de corps résiduel
$k$. Les hypothèses passent à $\bar C = C / \mathfrak{m} C$ et à la dérivation
réduite $\bar D$, et l'énoncé sur les invariants donne
$\operatorname{Ker}\bar D = k$. Un endomorphisme nilpotent d'un $k$-espace de
dimension $p$ dont le noyau est une droite a un seul bloc de Jordan, donc
$\bar D^{p-1} \neq 0$. Choisissons $x \in C$ avec $\bar D^{p-1} \bar x \neq
0$ : alors $D^{p-1} x$ est dans $\operatorname{Ker} D = A$ et a une image non
nulle dans $k$ ; c'est une unité $c$ de $A$, et $u = c^{-1} D^{p-2} x$
convient.\footnote{La page raisonne autrement, dans un passage encadré et
traversé de deux obliques, dont on ne sait s'il est annulé : $A$ étant local,
l'image de $D$ est libre de rang $p - 1$ (isomorphe à $C/A$), et elle
contient $A$ parce que tout $v = \sum \lambda_i D^i$ nul sur $D(C)$ vérifie
$vD = 0$, donc $\lambda_0 = \cdots = \lambda_{p-2} = 0$, donc $v$ est nul
sur $A$. Le calcul est juste, mais conclure $A \subset D(C)$ de
$D(C)^{\circ} \subset A^{\circ}$ demande, par le lemme, que $D(C)$ soit
facteur direct dans $C$ --- ce qui est vrai sur un corps et que la page
n'établit pas sur un anneau local. La réduction au corps résiduel ci-dessus
est la nôtre. Une note tournée dans la marge, presque illisible, invoque un
endomorphisme $\Theta$ avec $\Theta(1) = 0$ et $\Theta(u) = 1$ ; on ne la
reconstitue pas.}

\pagerange{20}{21}
\subsection*{Tout torseur est $A[U]/(U^p - a)$}

\emph{Si $Du = 1$, les $u^i$ ($0 \le i \le p-1$) forment une base de $C$ sur
$A$.} Posons $a = u^p$, qui est dans $A$, et $C_a = A[U]/(U^p - a)$ avec la
dérivation $D_a U = 1$. Par la Proposition, $C_a$ est un torseur. L'homomorphisme
$C_a \to C$, $U \mapsto u$, commute aux dérivations, donc est un morphisme de
torseurs sous $\alpha_p$ ; et un morphisme entre deux torseurs sous un même
groupe est un isomorphisme.\footnote{La page énonce ce dernier fait entre
crochets, pour deux pseudo-torseurs dans une catégorie $\mathcal{C}$
quelconque, sous une condition dont la lecture est très incertaine (le mot
souligné lu « monomis », un carré $G \times E \to G \times E'$ au-dessus de
$E \times E \to E' \times E'$) ; elle l'applique à la catégorie des anneaux
finis sur $A$ et se demande s'il faut $A$ noethérien. Pour de vrais torseurs
il n'en est rien : après un changement de base fidèlement plat qui trivialise
les deux, un morphisme équivariant $G \to G$ est une translation, donc un
isomorphisme, et cela descend. Pour de simples pseudo-torseurs l'énoncé est
faux (l'objet vide s'envoie dans tout pseudo-torseur).}

Le début de la page 21 contient une autre démonstration, barrée, de
l'indépendance linéaire des $u^i$ : d'une relation $\sum \lambda_i u^i = 0$,
on tire une contradiction en appliquant une puissance de $D$. Il faut
l'appliquer au \emph{dernier} coefficient non nul $\lambda_\beta$ : $D^\beta$
donne alors $\beta!\,\lambda_\beta = 0$, donc $\lambda_\beta =
0$.\footnote{La page prend le \emph{premier} coefficient non nul
$\lambda_\alpha$ et écrit que $D^\alpha$ donne $\alpha!\,\lambda_\alpha = 0$ ;
c'est faux, car les termes d'indice supérieur laissent des puissances de $u$.
Le passage est encadré et barré, peut-être pour cette raison ; la lecture
« $\alpha!\,\lambda_\alpha = 0$ » est d'ailleurs incertaine. Cet argument ne
donnerait de toute façon que l'indépendance, pas la génération.}

\pagerange{22}{23}
\subsection*{Le « Théorème » et le « Corollaire »}

\begin{quote}
\textbf{Théorème} (page 22). \emph{L'ensemble $H^1(A, \alpha_p)$ des classes
de torseurs sous $\alpha_p$ s'identifie au groupe additif $A/A^p$, où
$A^p = \{x^p : x \in A\}$. À $a \in A$ correspond la classe de
$C_a = A[U]/(U^p - a)$, munie de la dérivation $D_a U = 1$ ; et
$C_a \simeq C_b$ si et seulement si $a \equiv b \pmod{A^p}$. Tout torseur
$C$ est isomorphe à un $C_a$, avec $a = u^p$ pour un $u$ tel que $Du = 1$,
lui-même déterminé à l'addition près d'un élément de $A$.}
\end{quote}

Le dernier point est la section précédente. Pour le critère d'isomorphie :
un isomorphisme équivariant $C_a \to C_b$ envoie $U$ sur un élément $v$ tel
que $D_b v = 1$, donc $v = U + c$ avec $c \in A$, et $a = v^p = b + c^p$.
Comme $A$ est de caractéristique $p$, $A^p$ est un sous-anneau et le
quotient $A/A^p$ un groupe.\footnote{Ce qui précède vaut pour les torseurs à
algèbre libre, ceux que la page considère. Pour un torseur quelconque,
l'algèbre est localement libre de rang $p$, et l'argument de la section
précédente --- qui est local, puis recollé par
$H^1(\operatorname{Spec} A, \mathcal{O}) = 0$ --- fournit encore un $u$
global ; l'énoncé vaut donc pour tous les torseurs.}

La page ajoute : « On n'a pas prouvé que la loi additive naturelle de
$H^1(A; B)$ correspond à celle de $A/A^p$ ». Elle y correspond. La
bijection est le cobord $\delta \colon A \to H^1(A, \alpha_p)$ de la suite
exacte de Frobenius des pages 24 et 25 --- $\delta(a)$ est le torseur des
racines $p$-ièmes de $a$, c'est-à-dire $\operatorname{Spec} C_a$ --- et un
cobord est un homomorphisme de groupes.\footnote{La fin de la phrase de la
page est mal lue (« \ldots{} paraît être pas facile ») ; on ne sait si elle
juge la vérification facile ou difficile. Le recours au cobord est le nôtre ;
la page 27 se sert d'ailleurs de cette loi, sans y revenir.}

\begin{quote}
\textbf{Corollaire} (pages 22 et 23). \emph{Une $A$-algèbre commutative $C$
admet une structure de torseur sous $\alpha_p$ si et seulement si elle admet
une $p$-base réduite à un élément, c'est-à-dire un $u$ tel que les $u^i$
($0 \le i \le p-1$) forment une base de $C$ sur $A$ et que $u^p \in A$.\footnote{La page écrit « i.e. un $u \in C$ tel que les $u^i$ forment une
base de $C$ sur $A$ », sans répéter $u^p \in A$, que contenait la définition
de la page 17. La condition est nécessaire : sur un corps $k$ de
caractéristique $p$, l'algèbre $k[u]/(u^p - u)$ a la base $(u^i)$ mais est
étale, et n'est pas un torseur sous $\alpha_p$ ; aucune dérivation ne
vérifie $Du = 1$, car on aurait $0 = D(u^p - u) = -Du$.} Une
telle structure équivaut à la donnée d'une classe de tels $u$ modulo $A$, ou
encore de la dérivation $D$ telle que $Du = 1$.}
\end{quote}

Enfin, la page décrit les
sous-$A$-hyperalgèbres $\mathcal{U}' \subset \operatorname{End}_A(C)$,
isomorphes à $A[X]/(X^p)$, qui font de $C$ un torseur. Une telle
$\mathcal{U}'$ est $A[D]$, et elle détermine $D$ au produit près par une unité
de $A$, puisque les éléments primitifs de $A[X]/(X^p)$ sont les $cX$ et que
ses automorphismes sont les $X \mapsto cX$ avec $c \in A^*$. Remplacer $D$ par
$c^{-1} D$ remplace $u$ par $cu$ ; la donnée de $\mathcal{U}'$ équivaut donc à
celle d'une classe de $p$-bases modulo $u \mapsto au + b$, avec $a \in A^*$ et
$b \in A$, c'est-à-dire d'un point de $(C/A)/A^*$.\footnote{La page écrit que
ces données sont en correspondance biunivoque avec \emph{les} éléments de
l'espace projectif $\mathbf{P}_A(C/A) = (C/A)/A^*$. En général on n'obtient
qu'une partie de $(C/A)/A^*$ : les classes dont un représentant engendre une
$p$-base. Sur un corps $k$, quand $C$ est un corps radiciel de degré $p$, tout
$u \notin k$ convient et l'on obtient tout l'espace projectif privé de la
classe nulle, de dimension $p - 2$. Le passage de la page 23 est par endroits
illisible.} En réponse partielle à la première « Question » de la page 18 :
pour $C$ radiciel de degré $p$ sur un corps, les façons de présenter
$\operatorname{End}_k(C)$ comme $C \otimes \mathcal{U}'$ avec
$\mathcal{U}' \simeq A[X]/(X^p)$ forment un espace projectif.

\pagerange{24}{26}
\section*{Le point de vue global : la suite de Frobenius (pages 24 à 26)}

Sous le titre « Autre façon d'obtenir ces résultats, et des résultats
\emph{globaux} », la page reprend tout en une suite exacte. Soit $X$ un
préschéma sur $\mathbb{F}_p$ --- un schéma, dirait-on aujourd'hui. Le groupe
additif $\mathbb{G}_{a,X}$ a pour algèbre de Hopf $\mathcal{O}_X[T]$, avec
$T$ primitif ; le Frobenius $\Phi$, donné par $T \mapsto T^p$, est un
endomorphisme du groupe, et son noyau --- l'image inverse de la section unité,
d'algèbre $\mathcal{O}_X[T]/(T^p)$ --- est $\alpha_{p,X}$.\footnote{La page
écrit que l'idéal d'augmentation « est engendré par $T^p$ ». C'est $T$ qui
engendre $\operatorname{Ker}\varepsilon$, et $T^p = \Phi(T)$ qui engendre
l'idéal $\Phi(\operatorname{Ker}\varepsilon)$ ; la conclusion,
$\mathcal{O}_X[T]/(T^p)$, est juste. La page appelle « hyperalgèbre »
l'algèbre de Hopf $\mathcal{O}_X[T]$.} D'où la suite exacte
\[
0 \longrightarrow \alpha_{p,X} \longrightarrow \mathbb{G}_{a,X}
\xrightarrow{\ \Phi\ } \mathbb{G}_{a,X} \longrightarrow 0 .
\]
Qu'elle soit exacte à droite est le point délicat. La page affirme que
$\Phi$ est surjectif « en un sens technique suffisant » : tout morphisme
$\mathbb{G}_{a,X} \to Y$ invariant sous l'action de $\mathbb{G}_{a,X}$ par
$\Phi$ se factorise par $X$ --- autrement dit $\Phi$ est un épimorphisme
effectif.\footnote{En marge, un « ? » et, de biais, « justifier ? ».
L'argument esquissé par la page (une injectivité sur les algèbres, « du moins
si $Y$ est affine ») est en partie illisible. L'énoncé correct est que
$\Phi$ est fini et fidèlement plat, donc un épimorphisme de faisceaux pour la
topologie \emph{plate} (fppf). Il ne l'est pas pour la topologie de Zariski ni
pour la topologie étale : extraire une racine $p$-ième n'est pas possible
étale-localement. La cohomologie qui suit est donc la cohomologie plate.}
La suite longue de cohomologie, où $A = \Gamma(X, \mathcal{O}_X)$, s'écrit
\[
0 \to H^0(X, \alpha_p) \to A \xrightarrow{\ F\ } A \to H^1(X, \alpha_p)
\to H^1(X, \mathcal{O}_X) \xrightarrow{\ F\ } H^1(X, \mathcal{O}_X),
\]
en utilisant que la cohomologie plate de $\mathbb{G}_a$ est la cohomologie
ordinaire de $\mathcal{O}_X$.\footnote{Ce dernier point --- la comparaison
des cohomologies plate et de Zariski pour un faisceau quasi cohérent --- est
un théorème de Grothendieck (descente fidèlement plate) ; la page l'utilise
sans le dire. La phrase entre crochets qui suit la suite sur la page est en
partie illisible.} Donc $H^1(X, \alpha_p)$ est une extension
\[
0 \to A/A^p \to H^1(X, \alpha_p) \to \operatorname{Ker}\bigl(F \mid
H^1(X, \mathcal{O}_X)\bigr) \to 0 .
\]
Pour $X$ affine, $H^1(X, \mathcal{O}_X) = 0$ et l'on retrouve le Théorème de
la page 22. Pour $X$ une variété complète sur un corps algébriquement clos
$k$, $A = k$ est parfait, $A/A^p = 0$, et
\[
H^1(X, \alpha_p) = \operatorname{Ker}\bigl(F \mid H^1(X, \mathcal{O}_X)\bigr),
\]
le sous-espace du $k$-espace de dimension finie $H^1(X, \mathcal{O}_X)$ tué
par le Frobenius.\footnote{$F$ est $p$-linéaire ($F(\lambda x) = \lambda^p
F(x)$), de sorte que son noyau est bien un sous-espace vectoriel. Le mot
« annulé » est d'une lecture incertaine. L'égalité $A = k$ suppose $X$ connexe
et réduit, ce que contient le mot « variété ».}

\pagerange{26}{27}
\section*{Une remarque sur les revêtements de $k[[X]]$ (pages 26 et 27)}

Soit $A = k[[X]]$, $k$ algébriquement clos. Un torseur non trivial sous
$\alpha_p$ est $A[T]/(T^p - f)$ avec $f \notin A^p$, et l'on peut retrancher
de $f$ toute puissance $p$-ième, c'est-à-dire ses termes d'exposant divisible
par $p$. Soit $n$ le plus petit exposant non divisible par $p$ qui reste : en
changeant d'uniformisante ($Y = X\,\epsilon^{1/n}$ pour une unité $\epsilon$
convenable), l'anneau devient
\[
A[X^{n/p}] \simeq A[T]/(T^p - X^n), \qquad n \ge 1,\ p \nmid n ;
\]
pour $f \in A^p$ on obtient l'anneau $A[T]/(T^p)$.\footnote{La page borne $n$
par $p - 1$ (la borne inférieure est illisible). C'est trop restrictif :
$n$ parcourt tous les entiers positifs premiers à $p$, et les anneaux obtenus
sont deux à deux non isomorphes. Dans $k[[Z]]$ avec $Z^p = X$, l'anneau est
$k[[Z^p, Z^n]]$, dont le semi-groupe des valuations, engendré par $p$ et $n$,
est un invariant ; pour $n = p+1$ il ne contient aucun entier de $1$ à
$p - 1$, et ne provient d'aucun $n \le p-1$. Il s'agit ici de l'anneau, non
de la structure de torseur, que l'unité $\epsilon$ modifie.} Si $n = 1$ c'est
un anneau local régulier ; sinon il ne l'est pas. Pour $p = 3$ et $n = 2$, on
obtient $A[X^{2/3}]$, le sous-anneau de $k[[Y]] = k[[X^{1/3}]]$ formé des
séries sans terme en $Y$ --- le cusp.

Or, pour la loi de groupe de $H^1(A, \alpha_p) = A/A^p$, ce torseur non
régulier est la \emph{somme} de deux torseurs réguliers : ceux des racines
$p$-ièmes de $X$ et de $-X + X^2$, puisque
$X + (-X + X^2) = X^2$.\footnote{La page dit cela pour $p$ quelconque, après
l'exemple $p = 3$ ; l'addition est celle de $A/A^p$, dont la page 22 disait
n'avoir pas prouvé qu'elle est celle de $H^1$.} La régularité ne passe donc
pas à la somme. Elle ne passe même pas au produit fibré : si $C_1$ et $C_2$
sont ces deux anneaux, $C_1 \otimes_A C_2 \simeq C_1[T]/(T^p)$ a des
nilpotents. En effet, dans $C_1$, on a $X = s^p$ et
$-X + X^2 = (-s + s^2)^p$, de sorte que l'équation de $C_2$ devient une
puissance $p$-ième parfaite.\footnote{Ce calcul est le nôtre ; la page
affirme l'isomorphisme $C_1 \otimes_A C_2 \simeq C_1[T]/(T^p)$ sans le
justifier. La somme des deux torseurs est le quotient de ce produit fibré par
l'antidiagonale de $\alpha_p \times \alpha_p$.}

\end{document}
